% !htbp
\begin{table}[H]

  \centering
  % ============================= Outer container (single prompt)
  \begin{tcolorbox}[
    rounded corners,
    arc=3pt,
    enhanced,
    width=1\textwidth,
    colframe=codegreen, % <----- CHANGED
    colback=white,
    title=\textbf{Prompts for Generating Kernel Code in NPUKernelBench},
    fonttitle=\bfseries\small,
    boxrule=0.6pt,
    left=2pt,
    right=2pt,
    top=0pt,
    bottom=1pt
  ]
    \begin{tcolorbox}[
      sharp corners,
      colframe=reforange,
      colback=reforange!5,
      colbacktitle=reforange!15,
      coltitle=reforange,
      boxrule=0.4pt,
      title={\textbf{Instructions}},
      fonttitle=\bfseries\small,
      left=0.5pt,
      right=0.5pt,
      top=0.5pt,
      bottom=0.5pt
    ]
{\small Please see Table~\ref{tab:prompt_ins}.}
\end{tcolorbox}

\begin{tcolorbox}[
      sharp corners,
      colframe=taskblue,
      colback=taskblue!5,
      colbacktitle=taskblue!15,
      coltitle=taskblue,
      boxrule=0.4pt,
      title={\textbf{API\_Desc.md}},
      fonttitle=\bfseries\small,
      left=0pt,
      right=0pt,
      top=0pt,
      bottom=0pt
    ]
\begin{lstlisting}[style=mdstyle]
# aclnnBasicMatmul 
## Function Description 
### Kernel Function 
This Ascend C kernel performs the multiplication of two 2D matrices, A and B. It is a core component in deep learning models, such as linear layers and attention mechanisms. It takes two 2D tensors that satisfy the matrix multiplication rule and outputs their product.
### Computational Formula 
Assume the input tensors are $A$ ($m \times k$) and $B$ ($k \times n$). The output tensor $C$ ($m \times n$) is computed as: $$C_{ij} = \sum_{p=1}^{k} A_{ip} B_{pj}$$, where $i$ ranges from $1$ to $m$, and $j$ ranges from $1$ to $n$. $C_{ij}$ denotes the element in the $i$-th row and $j$-th column of $C$.
### Computation Process and Type Conversion
To maintain high numerical precision during large-scale accumulation and effectively prevent data overflow, this kernel adopts a high-precision accumulation strategy during computation. The process is as follows: 
1. The kernel receives two input tensors  `a` and `b`, both of data type `float16`.
2. During the multiply-accumulate computation, the internal accumulator uses the `float32` data type. In other words, the product results of `float16` inputs are first converted to `float32` before accumulation.
3. After all accumulation operations are completed, a result tensor of type `float32` is obtained.
4. Finally, the `float32` result tensor is converted back to `float16` as the final output.
## Interface Definition
### Kernel Prototype Definition Interface
#### Input 
- a: Device-side aclTensor corresponding to A in the formula; supports float16, 2D, and ND format.
- b: Device-side aclTensor corresponding to B in the formula; supports float16, 2D, and ND format.
#### Output 
- c: Device-side aclTensor corresponding to C in the formula; supports float16, 2D, and ND format.
#### Attr 
- None
## Constraints and Limitations
* The input tensors `a` and `b` currently support only the `float16` data type.
* The input tensors `a` and `b` must both be two-dimensional matrices.
* The second dimension (number of columns) of `a` must be equal to the first dimension (number of rows) of `b`.
* The input tensors support only the ND data format.
\end{lstlisting}
 \end{tcolorbox}
 
% --------------- Part 2: Previous CUDA implement with scores
\begin{tcolorbox}[
      sharp corners,
      colframe=refgreen,
      colback=refgreen!5,
      colbacktitle=refgreen!15,
      coltitle=refgreen,
      boxrule=0.4pt,
      title={\textbf{Host template}},
      fonttitle=\bfseries\small,
      left=0pt,
      right=0pt,
      top=-4pt,
      bottom=-4pt
    ]
      \begin{center}
        \begin{lstlisting}[style=mystyle, frame=none,framerule=0pt, language=Python]
#include "register/op_def_registry.h" 
#include "tiling/platform/platform_ascendc.h" 
namespace optiling { 
static ge::graphStatus TilingFunc(gert::TilingContext *context) 
{ 
    context->SetBlockDim(platform_ascendc::PlatformAscendCManager::GetInstance()->GetCoreNumAic()); 
    return ge::GRAPH_SUCCESS; 
} 
} // namespace optiling 
namespace ops { 
class BasicMatmul : public OpDef { 
public: 
    explicit BasicMatmul(const char *name) : OpDef(name) 
{ 
        this->Input("a") 
            .ParamType(REQUIRED) 
            .DataType({ge::DT_FLOAT16}) 
            .Format({ge::FORMAT_ND}); 
        this->Input("b") 
            .ParamType(REQUIRED) 
            .DataType({ge::DT_FLOAT16}) 
            .Format({ge::FORMAT_ND}); 
        this->Output("c") 
            .ParamType(REQUIRED) 
            .DataType({ge::DT_FLOAT16}) 
            .Format({ge::FORMAT_ND}); 
        this->AICore() 
            .SetTiling(optiling::TilingFunc) 
            .AddConfig("ascend910_93") 
            .AddConfig("ascend910b"); 
    } 
}; 
OP_ADD(BasicMatmul); 
} // namespace ops
\end{lstlisting}
\end{center}
\end{tcolorbox}

\begin{tcolorbox}[
      sharp corners,
      colframe=refgreen,
      colback=refgreen!5,
      colbacktitle=refgreen!15,
      coltitle=refgreen,
      boxrule=0.4pt,
      title={\textbf{Kernel template}},
      fonttitle=\bfseries\small,
      left=0pt,
      right=0pt,
      top=-4pt,
      bottom=-5pt
    ]
      \begin{center}
        \begin{lstlisting}[style=mystyle, frame=none,framerule=0pt, language=Python]
#include <kernel_operator.h> 
using namespace AscendC; 
extern "C" __global__ __aicore__ void basic_matmul(GM_ADDR a, GM_ADDR b, GM_ADDR c, GM_ADDR workspace, GM_ADDR tiling) 
{ 
}
\end{lstlisting}
\end{center}
\end{tcolorbox}
\end{tcolorbox}
\vspace{-10pt}
\caption{An example of the structured prompt used for Ascend kernel generation, consisting of three components: API description, host/kernel-side template.}
\label{tab:kernel_gen}
\end{table}


